% !TeX root = ../paper.tex
\hypertarget{libero-smart}{%
\section{LIBERO-SMART}\label{libero-smart}}

LIBERO-Pro, LIBERO-Plus, and RoboTwin-IF evaluate generalization across changes in task goals,
layouts, observations, and language requirements. Such variations are important, but robust task
completion does not by itself establish that a policy follows the precise scope and execution
conditions of a current instruction. For example, rephrasing ``pick up the cup and place it on the
plate'' preserves the intended grasp-and-place sequence, whereas ``only approach the cup'' requires
suppressing its later stages. LIBERO-SMART makes these distinctions explicit through Semantic,
Visual Prompt, and Control tasks and instruction-consistent trajectory contracts
(Figure~\ref{fig:overview}(d)).

\hypertarget{instruction-generalization-over-familiar-skills}{%
\subsection{Instruction Generalization over Familiar
Skills}\label{instruction-generalization-over-familiar-skills}}

The benchmark asks whether familiar manipulation capabilities can be invoked under new requirements,
rather than requiring a new skill for every command. A policy that has learned grasping, transport,
and placement should be able to select the requested phase, bind its target to the current
observation, or redirect its continuation. We therefore retain some LIBERO scene assets and basic
skills while evaluating held-out skill combinations, language conditions, object--target relations,
and layouts. This separates familiar motor operations from the new conditions governing their use.

Different commands alone cannot exclude solutions based on familiar object locations. Each task
therefore has 10 substantially different macro layouts and 5 local position and pose variants per
layout, preserving instruction meaning and physical feasibility while changing object configurations
and target locations. For Control, models start from the same intermediate physical state within
each instance, so differences in preceding-stage execution do not change the difficulty of the
corrective continuation.

Reaching a target also does not establish that the full instruction was followed: a policy may first
manipulate a forbidden object, execute steps in the wrong order, or immediately undo the intended
result. Task contracts therefore jointly assess target correctness, process constraints, and result
stability, including applicable prohibitions, step-order constraints, and release or holding
requirements. Together, tasks, initializations, and contracts test whether improved
instruction--action correspondence yields complete instruction execution. Training--evaluation
separation and detailed construction and scoring rules are provided in
Appendix~\ref{appendix:smart-benchmark}.

\hypertarget{three-task-suites}{%
\subsection{Three Task Suites}\label{three-task-suites}}

LIBERO-SMART comprises three independently scored suites, Semantic,
Visual Prompt (Vis), and Control, with 40 tasks per suite and 50
initializations per task, totaling 6,000 evaluation
instances. Both direct policies and agentic systems are evaluated using
the same trajectory-level task contracts.

\textbf{Semantic.} This suite evaluates fine-grained semantic following
through requirements on the manipulated object, action scope, and
termination condition. It covers equivalent
semantic rewrites, designated subtasks of a complete episode, standalone
atomic actions, and logical constraints on object exclusivity, spatial
relations, pose, and order. For example, when an instruction only
requests picking up an object, subsequently placing it in the original
task container violates the requirement.

\textbf{Visual Prompt.} This suite evaluates multimodal instruction
understanding by testing whether models can bind explicit image cues to
the current object or spatial target. Cues include circles
identifying objects to manipulate, circles defining placement regions,
points specifying target locations, and trajectories guiding motion.
They refer to the current object or location in each initialization and
remain consistent with the language instruction and target evaluation.

\textbf{Control.} Starting from an intermediate physical state, this
suite evaluates mid-execution corrective control by testing whether
models adjust subsequent behavior according to a new instruction. Target correction changes the
intended final result, for
example redirecting a held object to another target. Process re-planning
changes how subsequent actions are executed, for example holding
stably, reversing an operation, or retaining completed parts while
rearranging the remaining steps.

